% =========================================================
% INTRODUCCIÓN
% =========================================================

\section*{Introducción}
% Forzamos que la introducción aparezca en el índice general sin número
\addcontentsline{toc}{section}{Introducción}
\vspace{0.5cm}

El siguiente trabajo está compuesto por diferentes marcos, los cuales
nos ayudaran a analizar y proponer resoluciones a las problemáticas
encontradas dentro de una Empresa Farmacéutica con respecto a su sistema
financiero. Por ello comenzamos con.

El marco conceptual, constituye la base de la investigación, tiene como
primordial propósito analizar, describir y registrar el problema
detectado en la Distribuidora Farmacéutica. Se establece el planteamiento del problema,
los objetivos que se desean, se formulan las preguntas que ayudaran a la
investigación y las hipótesis de trabajo. Esta sección no solo demuestra
las debilidades del control operativo vigente, sino que valida las
propuesta de solución la cual consiste en implementar un sistema
informático.

El marco teórico se centra en la implementación de la ciencia la cual es
indispensable para comprender y corroborar las variables propuestas en
el sistema contable y tecnológico bajo estudio. En este capítulo se
recopilan los antecedentes analíticos, las teorías de gestión financiera
enfocadas en el control de flujos de efectivo, la administración de
fondos fijos, la liquidación de viáticos y la mitigación de riesgos
operativos en cuentas por pagar. Asimismo, se incorporan los fundamentos
de la ingeniería de software, arquitectura modular, y el marco legal de
la República de Guatemala.

El marco metodológico corresponde al apartado operativo, la elección del
tipo de investigación. En este espacio se define el tipo de estudio, el
enfoque adoptado supuestos, estimación de recursos, entre otros. también
se determinan los instrumentos de medición, asegurando que todo el
proceso de recolección de datos sea transparente y optimo.